%% kapa: exam=1 round=20 year=2009 period=1 form=A subject=civil_law subject_ko=민법 seq=11 number=11
%% source: Q-Net 감정평가사 기출문제 (cst00309) artlSeq=5024178 file=extracted/1008/5024178/2009년 제20회 감정평가사 1차시험 기출문제/1교시.hwp
%% conversion: kordoc 4.24.51 -> tools/gen1.py (auto); review=proofread-sonnet; flags=none
\begin{question}{11}
민법상 법인의 설립에 관한 설명으로 옳지 \uline{않은} 것은? (다툼이 있으면 판례에 의함)
\begin{choices}
  \item 사단법인의 설립행위는 요식행위이다.
  \item 재단법인은 공익을 목적으로 해서만 설립할 수 있다.
  \item 재단법인의 설립을 허가할 것인지 여부는 주무관청의 재량에 속한다.
  \item 합동행위설에 따르면 사단법인을 설립하는 자 가운데 한 사람의 설립행위가 행위무능력을 이유로 취소되더라도 다른 설립자의 설립행위에는 영향을 미치지 않는다.
  \item 재단법인설립을 위해 부동산을 출연한 경우에 재단법인이 소유권등기를 갖추지 못하였다면 제3자에게 소유권의 취득을 주장하지 못한다.
\end{choices}
\end{question}
\answer{2}
